%!TeX root = ../main.tex

\chapter{Conclusion}\label{chapter:conclusion}

This work set out to determine whether Cloudflare R2 can replace Amazon S3 as a primary object store without compromising performance, and to quantify the cost implications of making that switch. Three research questions structured the investigation: whether R2 can saturate high-bandwidth EC2 network interfaces, how tail latency evolves under high concurrency, and whether the zero-egress cost model justifies migration. This chapter summarises the answers, states the practical recommendations that follow from them, acknowledges the study's limitations, and outlines directions for future work.

\section{Summary of Findings}

\textbf{RQ1 --- Throughput parity.} The r8gd.4xlarge's 15~Gbps interface was saturated at 95\% utilisation with no observable throttling detected from R2 under the tested conditions. At higher bandwidths, the Python benchmarking client became the binding constraint: per-core throughput of 0.8--1.5~Gbps (depending on core count) capped aggregate throughput at 40~Gbps (36 vCPUs), 78~Gbps (64 vCPUs), and 114~Gbps (128 vCPUs), while CPU utilisation held at 90--99\%. The 114~Gbps peak is therefore a lower bound on R2's actual serving capacity---within the same order of magnitude as the approximately 100~Gbps figure AWS documents for S3, a similarity consistent with shared client-side constraints rather than informative about either system's backend capacity. These results establish a lower bound on R2 performance under large-object, single-client workloads and should not be interpreted as a full characterisation of system capacity. A client implemented in a compiled language would be required to probe R2's true backend ceiling.

\textbf{RQ2 --- Latency under concurrency.} R2's tail latency grows roughly proportionally with concurrency, consistent with queueing-theory predictions. On the c6in.32xlarge-2, tripling concurrency from C768 to C2304 raised P99 by 3.7$\times$ (4,182~ms to 15,337~ms). The system degraded gracefully throughout: 100\% post-retry success was maintained at 2,304 concurrent connections, with a measured retry rate below 0.15\%. At saturation on the 100~MB configuration, 62\% of mean request time was accounted for by RTT---time spent before transfer began---suggesting that queueing or congestion pressure in the request path, rather than data transfer itself, is the dominant latency component under high load; the precise source (backend processing, TCP congestion, or client-side scheduling) cannot be isolated from client-side measurements. Reducing request granularity from 100~MB to 50~MB chunks at equal aggregate throughput cut RTT at peak concurrency from 5,925~ms to 251~ms on the c6in.32xlarge, consistent with reduced queueing pressure somewhere in the request path. The paired r8gd.4xlarge configurations (same chunk$\times$pipeline product) corroborate this at a different scale: the 50~MB / p=6 configuration achieved P99 of 13,674~ms at steady-state C=192, versus 14,957~ms for the 100~MB / p=3 configuration at steady-state C=96---a worse tail despite lower concurrency, attributable to the higher per-request service time of larger chunks.

\textbf{RQ3 --- Cost-performance trade-off.} R2 delivers 53--77\% monthly savings over S3 for egress-heavy workloads, with the exact figure depending on the egress-to-storage ratio. The one-time migration cost is dominated (95--98\%) by S3 egress charges, and break-even is typically reached within one to four months at medium and large scales. For the large-object bulk transfer workloads benchmarked in this study, R2's measured throughput is within the same order of magnitude as S3 at single-instance scale, with the client as the bottleneck in both cases; the cost savings are therefore realisable without observed performance penalty for this workload class.

\section{Practical Recommendations}

The findings support a workload-dependent migration strategy rather than a blanket recommendation.

\textbf{Strong case for R2.} Large-object batch pipelines, data lake archival, ML training data distribution, and media delivery are all strong candidates. These workloads are characterised by large sequential reads, high egress volume, and tolerance for per-request latencies in the hundreds of milliseconds. R2 achieves throughput comparable to S3 for these patterns while eliminating the cost component that dominates at scale. Globally distributed access patterns benefit additionally from R2's anycast routing, which reduces P95 latency by approximately 38\% for multi-region access compared to a fixed S3 region~\cite{cloudflare-r2-faster-s3}.

\textbf{Weaker case for R2.} Latency-sensitive workloads serving small objects (below 1~MB) are better served by S3. Third-party benchmarks report S3 P50 of 22~ms for 1~KB objects~\cite{tigris-benchmark} against R2's 606~ms---a 28$\times$ difference that no cost saving compensates for in interactive or real-time applications. Similarly, workloads that depend on S3-specific features---Object Lock, Intelligent-Tiering, S3 Select, or deep integration with AWS services such as Athena and Lambda---face non-trivial adaptation costs that should be factored into the migration decision.

\textbf{Instance and configuration sizing.} For R2 read workloads from EC2, instances with 16--64 vCPUs represent the best cost-performance operating point with a Python client: the r8gd.4xlarge achieved full NIC saturation at 15~Gbps, while the c6in.16xlarge reached 78\% of its 100~Gbps interface with half the vCPU count of the 128-vCPU c6in.32xlarge. For write workloads or migration pipelines, the EC2 internet-gateway multi-flow cap (5~Gbps for instances with fewer than 32 vCPUs; 50\% of provisioned bandwidth for 32+ vCPUs) is a documented sizing constraint that should be accounted for explicitly; this cap was not observed to be enforced in the download experiments conducted here, but outbound traffic for upload-heavy workloads may be subject to it. Chunk sizes of 50~MB with pipeline depth 6 are preferable to 100~MB/p3 when latency matters: they achieve equivalent throughput with significantly lower RTT at high concurrency.

\section{Limitations}

The study has several constraints that bound the generalisability of its conclusions. The Python client became the bottleneck above 15~Gbps, so all throughput figures above that point are lower bounds on R2's capacity. The single-object workload---byte-range GETs from one 9~GB object---may not reflect R2's behaviour under production-scale multi-object access distributions. Only two AWS regions (eu-north-1 Stockholm and eu-central-1 Frankfurt) were tested; results may differ in regions where Cloudflare's network topology or R2 point-of-presence placement differs. The S3 baseline data was collected with an earlier framework version and is used only for order-of-magnitude comparison. Finally, the 100\% reported success rates reflect post-retry outcomes; applications that cannot tolerate automatic retries face a non-zero transient failure rate at high concurrency.

\section{Future Work}

Several directions would strengthen and extend the findings of this study.

\textbf{S3 baseline with the current framework.} The S3 reference data used in this study was collected with an earlier version of R2-bench and is used only for order-of-magnitude comparison. Re-running the full benchmark suite against S3 Standard using the current framework---identical instance types, chunk sizes, concurrency levels, and pipeline depths---would produce a controlled head-to-head comparison and allow definitive conclusions about which system is the binding constraint under matched conditions.

\textbf{Broader R2 characterisation for the same workload class.} The current study covers six instance types across two AWS regions. Systematically expanding coverage to additional instance types (particularly those in the 15--50~Gbps range), additional concurrency and chunk-size combinations, and additional AWS and Cloudflare regions would produce a more comprehensive characterisation of R2's large-object read throughput. This would also test whether the step-down behaviour observed on the r8gd.4xlarge long-run experiments generalises to other ``up to'' instances, and whether the per-core throughput ceiling is stable across different EC2 generations. Such a dataset would provide the empirical breadth needed to support publication.

\textbf{Compiled-language client.} Implementing the benchmarking client in a compiled language (Go, Rust, or C++) would remove the Python bottleneck and allow probing R2's actual backend throughput ceiling, potentially revealing whether R2 can sustain throughput significantly above the 114~Gbps lower bound measured here.

\textbf{Multi-object and small-object workloads.} Expanding the workload to cover multi-object access patterns---millions of small keys drawn from a realistic access distribution---would test R2's metadata index performance and bucket-level parallelism, and would complement the existing third-party benchmarks on small-object latency.

\textbf{Write throughput.} A direct comparison of upload throughput under the same framework would complete the picture for organisations evaluating R2 for data ingest pipelines, where the EC2 internet-gateway outbound cap is a documented constraint for upload-heavy workloads.

\textbf{Network-level diagnostics.} All measurements in this study are client-side only; latency components such as TCP congestion, R2 backend processing delay, and Cloudflare PoP routing cannot be isolated. Augmenting the framework with path-level measurements---traceroutes to individual Cloudflare PoPs, RTT baselines, and packet-level timing---would allow the source of elevated pre-transfer RTT under high concurrency to be attributed more precisely, and would characterise how performance varies with Cloudflare PoP placement across AWS regions in North America, Asia-Pacific, and other geographies.
